\section{Appendix A: Soundness Theorem -- Full Proof}\label{app:proof}

We restate \cref{thm:soundness} for self-containment, then give the
full proof. The proof is structured as the three-step argument
sketched in~\cref{sec:soundness}: termination, safety, coverage.

\paragraph{Restated theorem.}
Let~\(\Lambda\) be a layout satisfying invariants I1--I3.
Let~\(\mathcal{G}\) be the recovery operator of \cref{def:G},
satisfying I4--I5. Let~\(\mathcal{E}\) be a family of physical
perturbations such that, for every event~\(e \in \mathcal{E}\) and
every state~\(s\), the number of corrupted symbols in any single
RS subblock of~\(e(s)\) does not exceed the correction capacity
\(t = \lfloor (n-k)/2 \rfloor\) declared by~\(\Lambda\). Then for
every initial consistent state~\(s_0 \models\) and every
\(e \in \mathcal{E}\), there exists~\(K \le |\Lambda|\) such that
\(\mathcal{G}^{K}(e(s_0)) \models\) or
\(\mathcal{G}^{K}(e(s_0)) = \bot\).

\paragraph{Proof.}
Let~\(s_0 \models\) and~\(e \in \mathcal{E}\). Define
\(s^{(0)} = e(s_0)\) and~\(s^{(i+1)} = \mathcal{G}(s^{(i)})\) for
\(i \ge 0\), terminating the sequence at the first~\(K\) such that
\(s^{(K)} \models\) or~\(s^{(K)} = \bot\).

\textbf{Termination.}
Define~\(N(s)\) as in \cref{lem:descent}: the number of regions
of~\(\Lambda\) on which the syndrome of~\(s\) is non-zero.
By construction~\(N(s) \in \{0, 1, \dots, |\Lambda|\}\).
Furthermore, \(N(s) = 0\) iff~\(s \models\) (\cref{def:consistent},
clause~1 and~2 are equivalent to all syndromes being zero; clause~3
follows from~I3 once syndromes are zero on dual-layer regions;
clause~4 follows from~I2 since sentinels are part of the redundancy
coverage). By \cref{lem:descent}, for every~\(i\) such that
\(s^{(i)} \notin \{\bot\} \cup \{s : s \models\}\), we have
\(N(s^{(i+1)}) \le N(s^{(i)}) - 1\). The sequence~\((N(s^{(i)}))_i\)
is therefore strictly decreasing on the non-terminating part, and
bounded below by~0; it must reach~0 in at most~\(|\Lambda|\) steps,
or hit~\(\bot\) earlier. In either case~\(K \le |\Lambda|\).

\textbf{Safety.}
We show that the sequence cannot terminate in a state~\(s^{(K)}
\notin \{\bot\} \cup \{s : s \models\}\). Suppose for contradiction
that it does. By the termination clause, the sequence terminates
when~\(s^{(K)} \models\) or~\(s^{(K)} = \bot\), neither of which
holds by hypothesis: contradiction. The sequence therefore terminates
in one of the two safe states.

We further show that no intermediate state silently masks a
corruption. Inspection of \cref{def:G} shows that the only ways
the operator returns~\(s' \neq \bot\) on input~\(s\) are: (a) all
syndromes are zero (skip, return~\(s\) unchanged), or (b) RS
correction succeeded \emph{and} CRC re-verification succeeded
\emph{and} sentinels are at canonical values. Case~(a) preserves
consistency by~\cref{prop:idem}. Case~(b) restores
consistency on the affected region by I3 and I2. Any state
returned by the operator therefore either is consistent on the
affected region or is~\(\bot\); silent corruption is impossible
by construction.

\textbf{Coverage.}
The hypothesis on~\(\mathcal{E}\) asserts that each event~\(e\)
induces at most~\(t\) symbol errors per RS subblock. The
Reed--Solomon code, by its bounded-distance decoding property,
decodes correctly any received word within Hamming distance~\(t\)
of a code word, which is exactly the regime in which step~2 of
\cref{def:G} returns a valid corrected codeword. The CRC
re-verification of step~3 is, by I3, a strict tightening: any
corrected word that fails CRC is rejected, even if RS accepted it.
Therefore the operator, applied iteratively to~\(e(s_0)\), restores
each affected region to its consistent state, until either all
regions are consistent (\(N = 0\),~\(s^{(K)} \models\)) or a region
is found whose RS-corrected codeword fails CRC, indicating
corruption beyond the design assumption (out of~\(\mathcal{E}\)),
in which case~\(\bot\) is returned.

This concludes the proof. \(\mathchar"1003\)

\paragraph{What the proof does not show.}
The proof is for a single application chain on a fixed
perturbation~\(e\). It does not address (a) interleaved events
arriving during the recovery loop itself, which would require a
concurrent-execution model; (b) probabilistic statements about the
distribution of~\(e\) in~\(\mathcal{E}\), which would require a
measure on the perturbation family; (c) machine-checkability,
which would require mechanization in a proof assistant. All three
are noted as v3 work in \cref{sec:limits}.
